% 第 16 题频率分布直方图：柱高与刻度按固定源 PDF 矢量路径实测重绘（组距 0.5，共 9 组）。
% 要点：最左组 [0,0.5) 有柱（高 0.08，紧贴 y 轴，勿漏）；纵轴 7 个刻度 0.04/0.08/0.12/0.16/a/0.42/0.50（无 0 标签）；
% 水平虚线自纵轴向右、止于「柱高恰等于该刻度的最后一根柱」的左边缘，不贯穿全图；
% 横轴标题在轴线右上方（轴线越过 4.5 刻度后仍在延伸，标题在其上方），全部坐标以 \ux/\uy 统一换算。
\begin{tikzpicture}[>=stealth,every node/.style={font=\small}]
\def\ux{1.594}   % 横轴：1 吨（0.5 吨 = 0.797 cm）
\def\uy{9.96}    % 纵轴：1 单位（0.04 = 0.3984 cm）
% 坐标轴（横轴越过末刻度延伸至 5.5 单位，与原图一致；纵轴至 0.60 单位）
\draw[->] (0,0)--(5.50*\ux,0);
\node[anchor=south west,inner sep=1pt] at (4.686*\ux,0.015*\uy) {月均用水量（吨）};
\draw[->] (0,0)--(0,0.60*\uy);
\node[anchor=south west] at (0.115*\ux,0.558*\uy) {\begin{tabular}{c}频率\\\hline 组距\end{tabular}};
% 纵轴刻度标签（含字母 a，无 0）
\foreach \v in {0.04,0.08,0.12,0.16,0.42,0.50} {
  \node[left] at (-0.06,\v*\uy) {\(\v\)};
}
\node[left] at (-0.06,0.30*\uy) {\(a\)};
% 横轴刻度标签
\foreach \t in {0,0.5,1,1.5,2,2.5,3,3.5,4,4.5} {
  \node[below] at (\t*\ux,-0.03) {\(\t\)};
}
% 水平虚线（止于指定柱左缘：0.50→2.0，0.42→1.5，a(0.30)→2.5，0.16→0.5，0.12→3.0，0.08→3.5，0.04→4.0）
\draw[dashed] (0,0.50*\uy)--(2.0*\ux,0.50*\uy);
\draw[dashed] (0,0.42*\uy)--(1.5*\ux,0.42*\uy);
\draw[dashed] (0,0.30*\uy)--(2.5*\ux,0.30*\uy);
\draw[dashed] (0,0.16*\uy)--(0.5*\ux,0.16*\uy);
\draw[dashed] (0,0.12*\uy)--(3.0*\ux,0.12*\uy);
\draw[dashed] (0,0.08*\uy)--(3.5*\ux,0.08*\uy);
\draw[dashed] (0,0.04*\uy)--(4.0*\ux,0.04*\uy);
% 柱（空心矩形，相邻柱共用竖边；柱高依次 0.08/0.16/0.30/0.42/0.50/0.30/0.12/0.08/0.04）
\foreach \i/\h in {0/0.08, 1/0.16, 2/0.30, 3/0.42, 4/0.50, 5/0.30, 6/0.12, 7/0.08, 8/0.04} {
  \draw ({0.5*\i*\ux},0) rectangle ({0.5*(\i+1)*\ux},{\h*\uy});
}
\end{tikzpicture}
